\documentclass[11pt]{article}

\usepackage[margin=1in]{geometry}
\usepackage{amsmath}
\usepackage{amssymb}
\usepackage{enumitem}

\setlength{\parindent}{0pt}
\setlength{\parskip}{0.55em}
\newcommand{\answerspace}[1]{\vspace{#1}}

\begin{document}

\begin{center}
{\Large \textbf{Week 5: Perceptron Learning with Binary Cross-Entropy}}\\
CSCE 5218 -- Deep Learning
\end{center}

\textbf{Name: Zack Butz} \hfill \textbf{Date: \today} 

\section*{Given}

For training example $k$, use
\[
i_k = \sum_{i=0}^{n} w_i x_{ki} = \mathbf{w}^{T}\mathbf{x}_k,
\qquad
y_k = \sigma(i_k) = \frac{1}{1+e^{-i_k}},
\]
where $t_k \in \{0,1\}$, $x_{k0}=1$, and
\[
L_k = -\left[t_k\ln(y_k)+(1-t_k)\ln(1-y_k)\right].
\]

\section*{1. Loss derivative}

\textbf{Question.} Derive $\partial L_k/\partial y_k$ and simplify it to one fraction.

\begin{enumerate}[label=\arabic*.]
    \item Expand the negative sign and differentiate each logarithm term with respect to $y_k$.
    \[
    L_k = -t_k\ln(y_k) - (1-t_k)\ln(1-y_k)
    \]
    \[
    \frac{\partial L_k}{\partial y_k}
    = -\frac{t_k}{y_k} + \frac{1-t_k}{1-y_k}
    \]

    \item Put the terms over the common denominator $y_k(1-y_k)$.
    \[
    \frac{\partial L_k}{\partial y_k}
    = -\frac{t_k(1-y_k)}{y_k(1-y_k)}
      + \frac{(1-t_k)y_k}{y_k(1-y_k)}
    \]
    \[
    \frac{\partial L_k}{\partial y_k}
    = \frac{-t_k(1-y_k) + (1-t_k)y_k}{y_k(1-y_k)}
    \]

    \item Show how the numerator simplifies.
    \[
    -t_k(1-y_k) + (1-t_k)y_k
    = -t_k + t_ky_k + y_k - t_ky_k
    = y_k-t_k
    \]

    \item Substitute the simplified numerator back into the fraction.
    \[
    \boxed{
    \frac{\partial L_k}{\partial y_k}
    = \frac{y_k-t_k}{y_k(1-y_k)}
    }
    \]
\end{enumerate}

\section*{2. Parameter gradients}

\textbf{Question.} Use the chain rule to derive $\partial L_k/\partial w_j$ for $j=0,\ldots,n$. Show every local derivative.

Use the computational path $w_j \rightarrow i_k \rightarrow y_k \rightarrow L_k$.

\begin{enumerate}[label=\arabic*.]
    \item Write the chain rule.
    \[
    \frac{\partial L_k}{\partial w_j}
    = \frac{\partial L_k}{\partial y_k}
      \frac{\partial y_k}{\partial i_k}
      \frac{\partial i_k}{\partial w_j}
    \]
    \item Fill in the three local derivatives.
    \[
    \frac{\partial L_k}{\partial y_k}=\frac{y_k-t_k}{y_k(1-y_k)} \qquad
    \frac{\partial y_k}{\partial i_k}=y_k(1-y_k) \qquad
    \frac{\partial i_k}{\partial w_j}=x_{kj}
    \]
    \item Substitute them and explicitly cancel common factors.
    \[
    \frac{\partial L_k}{\partial w_j}
    =
    \left(\frac{y_k-t_k}{y_k(1-y_k)}\right)
    \left(y_k(1-y_k)\right)
    x_{kj}
    \]
    \[
    \frac{\partial L_k}{\partial w_j}
    =
    (y_k-t_k)x_{kj}.
    \]
    \item Explain why the same formula includes the bias. What happens when $j=0$ and $x_{k0}=1$?
    The formula includes the bias so that the model can classify data which is not seperable \\
    through the origin acting like a boundary condition for the formula. When $j=0$ and $x_{k0}=1$, \\ 
    the gradient is simply $y_k-t_k$.
\end{enumerate}

\section*{3. Gradient-descent update}

\textbf{Question.} Write one update rule that applies to every $w_j$, including $w_0$, using learning rate $\eta$.

\begin{enumerate}[label=\arabic*.]
    \item Start from gradient descent:
    \[
    w_j \leftarrow w_j - \eta\frac{\partial L_k}{\partial w_j}.
    \]
    \item Substitute your simplified gradient.
    \[
    w_j \leftarrow w_j - \eta(y_k-t_k)x_{kj}
    \]
    \item When $x_{kj}>0$, explain what happens to $w_j$ when $t_k-y_k>0$, and when $t_k-y_k<0$. \\
    When $x_{kj}>0$ and $t_k-y_k>0$, the update is positive, so
    $w_j$ increases. When $t_k-y_k<0$, the update is negative, so
    $w_j$ decreases.
\end{enumerate}
\pagebreak
\section*{4. One numerical update}

Use
\[
\mathbf{x}_k=(1,1,2),\qquad
\mathbf{w}=(0,0.2,-0.1),\qquad
t_k=1,\qquad
\eta=0.1.
\]
Report nonlinear values to four decimal places.

\subsection*{4.1 Forward pass before the update}
\[
i_k = w_0x_{k0}+w_1x_{k1}+w_2x_{k2} = \mathbf{w}^T\mathbf{x}_k
\]
\[
y_k=\sigma(i_k)=\frac{1}{1+e^{-i_k}}
\]
\[
L_k=-\left[t_k\ln(y_k)+(1-t_k)\ln(1-y_k)\right]
\]
\[
\frac{\partial L_k}{\partial w_j}
=
(y_k-t_k)x_{kj}.
\]
\[
w_j \leftarrow w_j - \eta(y_k-t_k)x_{kj}
\]

\subsection*{4.2 Gradients}
\[
\frac{\partial L_k}{\partial w_0}=(\frac{1}{1+e^{0}}-1)1\qquad
\frac{\partial L_k}{\partial w_1}=(\frac{1}{1+e^{0}}-1)1\qquad
\frac{\partial L_k}{\partial w_2}=(\frac{1}{1+e^{0}}-1)2
\]
\[
\frac{\partial L_k}{\partial w_0}=-\frac{1}{2}\qquad
\frac{\partial L_k}{\partial w_1}=-\frac{1}{2}\qquad
\frac{\partial L_k}{\partial w_2}=-1
\]

\subsection*{4.3 Updated parameters}
\[
w_0^{\text{new}}=0 - 0.1(-0.5) = 0.05 \qquad
w_1^{\text{new}}=0.2 - 0.1(-0.5) = 0.25 \qquad
w_2^{\text{new}}=-0.1 - 0.1(-1) = 0  
\]

\subsection*{4.4 Forward pass after the update}
\[
i_k^{\text{new}}=(0.05)(1)+(0.25)(1)+(0)(2)=0.3
\]
\[
y_k^{\text{new}}=\sigma(0.3)=\frac{1}{1+e^{-0.3}}\approx 0.5744
\]
\[
L_k^{\text{old}}=-\left[1\ln(0.5)+(1-1)\ln(1-0.5)\right] = -\ln(0.5) \approx 0.6931
\]
\[
L_k^{\text{new}}=-\left[1\ln(0.5744)+(1-1)\ln(1-0.5744)\right] = -\ln(0.5744) \approx 0.5544
\]


Briefly verify whether the loss moved in the expected direction.
The loss decreased from $L_k\approx 0.6931$ to $L_k^{\text{new}}\approx 0.5544$, 
which is the expected direction since the value of $y_k$ is closer to its expected 
value of 1 after the update.

\section*{5. Comparison with squared error}

\textbf{Question.} For a sigmoid neuron, why does squared error retain $y_k(1-y_k)$ in its weight gradient while BCE does not? What can happen near sigmoid outputs of 0 or 1?

\begin{enumerate}[label=\arabic*.]
    \item If $E_k=\tfrac12(t_k-y_k)^2$, write its chain-rule gradient with respect to $w_j$.
    \[
    \frac{\partial E_k}{\partial w_j}= \frac{\partial E_k}{\partial y_k}
      \frac{\partial y_k}{\partial i_k}
      \frac{\partial i_k}{\partial w_j}
    \]
    \[
    \frac{\partial E_k}{\partial y_k} = -(t_k-y_k) = y_k-t_k
    \]
    \[
    \frac{\partial y_k}{\partial i_k} = y_k(1-y_k)
    \]
    \[
    \frac{\partial i_k}{\partial w_j} = x_{kj}
    \]
    \[
    \frac{\partial E_k}{\partial w_j} = (y_k-t_k)y_k(1-y_k)x_{kj}
    \]
    \item Explain one training consequence when $y_k$ is close to 0 or 1. \\
    When $y_k$ is close to 0 or 1, the term $y_k(1-y_k)$ becomes very small, 
    which can lead to attenuated gradients. This means that the weight updates
    will be very small, slowing down the learning process and making it 
    difficult for the model to converge to an optimal solution.
    This is one of the reasons why BCE is preferred over squared error 
    for classification tasks, as it avoids this issue and provides more stable 
    gradients due to removing the extra sigmoid derivative factor seen in section 2 part 3.

\end{enumerate}

\section*{AI assistance disclosure}

Initial bootstrapping for this homework was provided by ChatGPT. 
All subsequent work, including derivations, calculations, and explanations, was completed by me.

\end{document}
